\chapter{Multi-timescale kernels}
\label{ch:multitimescale}

\begin{motivation}
Real excitation often has \emph{several} timescales: a fast burst of immediate
reactions and a slow tail of long-term influence. A sum of exponentials captures
this, stays inside our closed-form machinery, and --- with a fixed bank of decay
rates --- is well-conditioned to fit.
\end{motivation}

\section{Sum-of-exponentials kernels}

Replace the single exponential by
\begin{equation}
\phi(t)=\sum_{q=1}^{Q}\kappa_q\theta_q e^{-\theta_q t},\qquad
\text{branching ratio }=\sum_q\kappa_q .
\end{equation}
Every result of Ch.~\ref{ch:solving} still holds per component, and the state-space
ODE (Ch.~\ref{ch:solvers}) simply carries $Q$ states instead of one.

\section{A closed-form compensator via the matrix exponential}

For a constant baseline the compensator has an exact closed form. The state vector
$u\in\R^Q$ obeys $\dot u=Au+a\mu$ with $A=a\mathbf 1^\top-\diag(b)$,
$a_q=\kappa_q\theta_q$, $b_q=\theta_q$, and $\xi=\mu+\mathbf 1^\top u$. Diagonalising
$A=V\Lambda V^{-1}$ gives, with $w=A^{-1}a$ and $p_i=(\mathbf 1^\top V)_i(V^{-1}w)_i$,
\begin{equation}
\Xi(t)=\mu\Big(t+\sum_i p_i\big(\tfrac{e^{\lambda_i t}-1}{\lambda_i}-t\big)\Big).
\end{equation}

\begin{whybox}[: why this matters for speed]
The pure-Python ODE-in-a-loop is too slow to call thousands of times inside an
optimiser. This matrix-exponential formula evaluates the compensator at all bucket
endpoints in $O(Q^3+mQ)$ with \emph{no} time grid --- turning a multi-minute fit
into a fraction of a second (Ch.~\ref{ch:neural} discusses the general remedy).
$A$ is invertible because $\sum_q\kappa_q<1$ makes it Hurwitz.
\end{whybox}

\section{Fitting: a fixed bank of timescales}

\begin{whybox}[: fix the decays, fit the weights]
Estimating both the rates $\theta_q$ and weights $\kappa_q$ is ill-posed
(timescales are weakly identified --- Ch.~\ref{ch:identif}, now $\times Q$). The
standard, well-conditioned remedy is to fix a \emph{bank} of decay rates (fast,
medium, slow) and fit only the non-negative weights $\kappa_q\ge0$, optionally with
an $\ell_1$ penalty to select a sparse set of active timescales.
\end{whybox}

\begin{lstlisting}[caption={Multi-timescale fit over a fixed decay bank.}]
from hawkes_calibration import fit_mbpp_ic_sumexp
bank = np.array([0.4, 1.0, 2.0, 5.0])     # fixed decay rates (slow ... fast)
res  = fit_mbpp_ic_sumexp(obs, counts, bank, l1=0.02)   # fit non-negative weights
res.kappas, res.thetas                     # weight per timescale; total = sum(kappas)
\end{lstlisting}

\begin{examplebox}[: total vs.\ split]
With a true 2-timescale kernel and a 3-rate bank, the fit recovers the \emph{total}
branching ratio well ($0.58$ vs.\ $0.60$) but spreads the weights across the bank ---
the per-timescale split is weakly identified, exactly as the $\kappa$--$\theta$
analysis predicts. Report the total and the rough fast/slow balance, not a precise
decomposition.
\end{examplebox}

\begin{recap}
Sum-of-exponentials kernels model fast+slow excitation, stay closed-form (matrix
exponential for the compensator), and are fit by fixing a decay bank and estimating
non-negative weights (with $\ell_1$). The total branching ratio identifies well;
the timescale split does not.
\end{recap}

\exercise{Show that a single-exponential kernel is the $Q=1$ case and recovers
Ch.~\ref{ch:solving}. Why does the matrix $A$ have all eigenvalues with negative
real part when $\sum_q\kappa_q<1$?}
